
\section{Optimization from Membership via Sampling}

In this chapter, we assumed the oracle can be computed exactly. However,
it is still open to determine the best runtime for reductions between
\emph{noisy} oracles, where the oracle provides answers to within
some bounded error. In particular, the question about minimizing (approximately)
convex functions under noisy oracles is an active research area, and
the gaps between the lower bound and upper bound for many problems
are still quite large (\cite{belloni2015escaping,Feldman2015statistical,bubeck2016kernel}
based on \cite{Kalai2006simulated}). We will see these methods in
detail later. For now, to put them in the context of oracles, we just
discuss the following theorem.
\begin{thm}[OPT via sampling]
\label{thm:OPTviaSampling} Let $F(x)=e^{-\alpha c^{T}x}1_{K}(x)$
for some convex set $K$ in $\R^{n}$ and vector $c\in\R^{n}.$ Suppose
$\min_{x\in K}c^{T}x$ is bounded. Let x be a random sample from the
distribution with density proportional to $F(x)$. Then, 
\[
\E\left(c^{T}x\right)\le\min_{K}c^{T}x+\frac{n}{\alpha}.
\]
\end{thm}

\begin{proof}
We will show that the worst case is when the convex set $K$ is an
infinite cone. WLOG\footnote{Without loss of generality.} assume
that the minimum is at $x=0.$ Replace every cross-section of $K$
along $c$ with an $(n-1)$-dimensional ball of the same volume as
the cross-section. This does not affect the expectation. Suppose the
expectation is $\E(c^{T}x)=a$. Next, replace the subset of $K$ with
$c^{T}x\le a$ with a cone whose base is the cross-section at $a$
and apex is at zero. This only makes the expectation larger. Next,
replace the subset of $K$ on the right of $a$ with the infinite
conical extension of the cone to the left of $a$. Again, this expectation
can only be higher. 

Now for this infinite cone, we compute, using $y=c^{T}x$, 
\begin{align*}
\E\left(c^{T}x\right) & =\frac{\int_{0}^{\infty}ye^{-\alpha y}y^{n-1}\,dy}{\int_{0}^{\infty}e^{-\alpha y}y^{n-1}\,dy}\\
 & =\frac{1}{\alpha}\frac{\int_{0}^{\infty}e^{-y}y^{n}\,dy}{\int_{0}^{\infty}e^{-y}y^{n-1}\,dy}\\
 & =\frac{1}{\alpha}\frac{n!}{(n-1)!}=\frac{n}{\alpha}\,,
\end{align*}
where the equality in the last line follows from properties of the
Gamma functions.
\end{proof}
The theorem says that if we sample according to $e^{-\alpha c^{T}x}$
for $\alpha=n/\varepsilon$, we will get an $\varepsilon$-approximation
to the optimum. However, sampling from such a density is not trivial.
Instead, we will have to go through a sequence of overlapping distributions,
starting with one that is easy to sample and ending with a distribution
that is focused close to the minimum. This method is known as \emph{simulated
annealing }and is the subject of Chapter \ref{chap:Annealing}. The
complexity of sampling is polynomial in the dimension and logarithmic
in a suitable notion of probabilistic distance between the starting
distribution and the target distribution. The sampling algorithm only
uses a membership (EVAL) oracle.
\begin{xca}
Extend Theorem \ref{thm:OPTviaSampling} by replacing $c^{T}x$ with
any convex function $f(x)$.
\end{xca}

\begin{problem*}
Given an approximately convex function $F$ on unit ball such that
$\max_{\norm x_{2}\leq1}|f(x)-F(x)|\leq\varepsilon/n$ for some convex
function $f$, how efficiently can we find $x$ in the unit ball such
that $F(x)\leq\min_{\norm x_{2}\leq1}F(x)+O(\varepsilon)$? The current
fastest algorithm takes $O(n^{4}\log^{O(1)}(n/\varepsilon))$ calls
to the noisy EVAL oracle for $F$ .
\end{problem*}

